% arXiv: [cs.DS]

\documentclass[11pt,letterpaper]{article}
\usepackage[margin=1in]{geometry}
\usepackage{amsmath,amssymb,amsthm}
\usepackage{graphicx}
\usepackage{pgfplots}
\pgfplotsset{compat=1.18}
\usepackage{algorithm}
\usepackage{algpseudocode}
 \usepackage{booktabs}
\usepackage{hyperref}
\usepackage{caption}

\newtheorem{theorem}{Theorem}[section]
\newtheorem{lemma}[theorem]{Lemma}
\newtheorem{corollary}[theorem]{Corollary}
\newtheorem{definition}[theorem]{Definition}
\newcommand{\keywords}[1]{\textbf{Keywords:} #1}
\title{Occupancy-Corrected Sparse Representation\\ for Cardinality Estimation}
\author{Luca Cappelletti\thanks{\texttt{luca.cappelletti@unifr.ch}, \texttt{https://orcid.org/0000-0002-1269-2038}}\\
  University of Fribourg}

\begin{document}
\maketitle
\keywords{cardinality estimation, HyperLogLog, streaming algorithms, sparse representation, hash list, occupancy correction, space-efficient counting}

\begin{filecontents*}{hll_plot_data.csv}
n,uncorrected,hllpp,hll_mle,corrected,setsketch_b2,setsketch_b1p2,setsketch_b1p001,bits
1,0.0010,0.0489,0.0481,0.0010,69.2214,69.2218,69.2218,0.000
2,0.0010,0.0978,0.0974,0.0010,55.3838,55.3839,55.3839,0.000
3,0.0010,0.2762,0.2748,0.0010,46.7725,46.7723,46.7723,0.000
4,0.0010,0.3887,0.3908,0.0010,40.6105,40.6105,40.6105,0.000
5,0.0010,0.4746,0.4781,0.0010,35.9880,35.9879,35.9879,0.000
6,0.0010,0.4836,0.4838,0.0010,33.1430,33.1431,33.1431,0.000
7,0.0010,0.5573,0.5567,0.0010,30.4137,30.4138,30.4138,0.000
8,0.0010,0.5772,0.5753,0.0010,28.6206,28.6205,28.6205,0.000
9,0.0010,0.7645,0.7641,0.0010,26.3874,26.3872,26.3872,0.000
10,0.0010,0.8474,0.8461,0.0010,25.2254,25.2254,25.2254,0.000
11,0.0010,0.8892,0.8875,0.0010,24.8927,24.8928,24.8928,0.000
12,0.0010,0.9869,0.9841,0.0010,23.9263,23.9267,23.9267,0.000
13,0.0010,1.0469,1.0427,0.0010,23.2700,23.2704,23.2704,0.000
14,0.0010,1.1247,1.1179,0.0010,22.0487,22.0489,22.0489,0.000
15,0.0010,1.1937,1.1845,0.0010,20.6809,20.6809,20.6809,0.000
16,0.0010,1.2425,1.2283,0.0010,19.4269,19.4265,19.4265,0.000
17,0.0010,1.3563,1.3408,0.0010,19.6434,19.6430,19.6430,0.000
18,0.0010,1.3874,1.3681,0.0010,19.4846,19.4842,19.4843,0.000
19,0.0010,1.4645,1.4474,0.0010,18.9703,18.9701,18.9701,0.000
20,0.0010,1.5197,1.5041,0.0010,18.5530,18.5529,18.5530,0.000
21,0.0010,1.5716,1.5544,0.0010,17.9406,17.9407,17.9408,0.000
22,0.0010,1.6219,1.6051,0.0010,17.3947,17.3950,17.3951,0.000
23,0.0010,1.6676,1.6543,0.0010,17.0896,17.0901,17.0902,0.000
24,0.0010,1.7062,1.6938,0.0010,16.6198,16.6204,16.6205,0.000
25,0.0010,1.7447,1.7295,0.0010,16.1787,16.1794,16.1794,0.000
26,0.0010,1.8006,1.7856,0.0010,15.7115,15.7120,15.7120,0.000
27,0.0010,1.8677,1.8524,0.0010,15.6851,15.6855,15.6855,0.000
28,0.0010,1.8746,1.8596,0.0010,15.3086,15.3089,15.3089,0.000
29,0.0010,1.8808,1.8672,0.0010,15.1595,15.1596,15.1596,0.000
30,0.0010,1.9240,1.9094,0.0010,14.8866,14.8865,14.8865,0.000
31,0.0010,1.9316,1.9221,0.0010,14.5503,14.5501,14.5501,0.000
33,0.0010,1.9780,1.9691,0.0010,13.7966,13.7962,13.7962,0.000
35,0.0010,1.9308,1.9236,0.0010,13.1014,13.1009,13.1009,0.000
37,0.0010,1.8766,1.8675,0.0010,12.6485,12.6480,12.6480,0.000
39,0.0010,1.8810,1.8754,0.0010,12.0491,12.0487,12.0488,0.000
41,0.0010,1.8410,1.8402,0.0010,11.9233,11.9234,11.9235,0.000
44,0.0010,1.7118,1.7106,0.0010,11.6600,11.6602,11.6602,0.000
47,0.0010,1.6346,1.6413,0.0010,11.4520,11.4525,11.4525,0.000
50,0.0010,1.7448,1.7221,0.0010,11.0975,11.0982,11.0983,0.000
53,0.0010,1.8191,1.8076,0.0010,10.6418,10.6424,10.6425,24.000
56,0.0010,1.8359,1.8260,0.0010,10.1721,10.1723,10.1724,24.000
59,0.0010,1.8143,1.8022,0.0010,10.1271,10.1268,10.1269,24.000
62,0.0010,1.7958,1.7796,0.0010,9.8494,9.8490,9.8490,24.000
66,0.0010,1.7806,1.7745,0.0010,9.5870,9.5863,9.5864,24.000
70,0.0010,1.8105,1.7817,0.0010,9.6490,9.6484,9.6486,24.000
74,0.0010,1.8015,1.7667,0.0010,9.5504,9.5498,9.5499,24.000
78,0.0010,1.7461,1.7303,0.0010,9.3930,9.3926,9.3929,24.000
82,0.0010,1.7641,1.7318,0.0010,9.2807,9.2803,9.2806,24.000
87,0.0010,1.8199,1.7697,0.0010,8.8940,8.8938,8.8940,24.000
92,0.0010,1.7612,1.7665,0.0010,8.5836,8.5833,8.5833,24.000
97,0.0010,1.8429,1.7936,0.0010,8.3641,8.3639,8.3640,24.000
102,0.0010,1.8095,1.7782,0.0010,8.1934,8.1933,8.1934,24.000
108,0.0010,1.8850,1.8297,0.0010,8.0063,8.0062,8.0063,24.000
114,0.0010,1.8626,1.8252,0.0010,7.9534,7.9527,7.9529,24.000
120,0.0010,1.8451,1.7911,0.0010,7.7304,7.7291,7.7292,24.000
126,0.0010,1.8745,1.8135,0.0010,7.5651,7.5641,7.5642,24.000
133,0.0010,1.8566,1.8000,0.0010,7.4396,7.4384,7.4386,24.000
140,0.0010,1.8339,1.8056,0.0010,7.1812,7.1800,7.1801,24.000
147,0.0027,1.8404,1.8344,0.0035,6.9719,6.9712,6.9714,24.000
155,0.0025,1.8292,1.8221,0.0034,6.8530,6.8509,6.8513,24.000
163,0.0024,1.8537,1.8311,0.0033,6.7307,6.7290,6.7293,24.000
172,0.0023,1.8529,1.8199,0.0033,6.6770,6.6774,6.6772,24.000
181,0.0022,1.8400,1.7976,0.0032,6.6600,6.6605,6.6604,24.000
191,0.0020,1.9099,1.8586,0.0032,6.5021,6.5023,6.5024,24.000
201,0.0019,1.8384,1.7991,0.0031,6.4483,6.4496,6.4497,24.000
212,0.0018,1.9257,1.8696,0.0031,6.4152,6.4170,6.4172,24.000
223,0.0018,1.9321,1.8732,0.0031,6.2323,6.2338,6.2342,24.000
235,0.0017,1.9339,1.8943,0.0030,6.2003,6.2008,6.2006,24.000
247,0.0032,1.9209,1.8634,0.0046,6.0810,6.0818,6.0814,24.000
260,0.0045,1.8970,1.8266,0.0075,5.9718,5.9723,5.9716,23.000
273,0.0057,1.9475,1.8803,0.0088,5.6698,5.6700,5.6696,23.000
287,0.0054,1.9720,1.9208,0.0087,5.5903,5.5894,5.5887,23.000
302,0.0052,1.9978,1.9469,0.0086,5.5632,5.5623,5.5618,23.000
318,0.0061,1.9835,1.9555,0.0097,5.3669,5.3672,5.3664,23.000
334,0.0070,1.9304,1.9043,0.0108,5.3343,5.3343,5.3341,23.000
351,0.0156,1.9475,1.8932,0.0230,5.2419,5.2431,5.2430,22.000
369,0.0286,1.9978,1.9148,0.0424,5.0442,5.0516,5.0515,21.000
388,0.0473,1.9700,1.9269,0.0683,4.8327,4.8405,4.8403,20.250
408,0.0574,1.9693,1.8748,0.0792,4.8214,4.8295,4.8299,20.000
429,0.0965,1.9767,1.8838,0.1182,4.7148,4.7251,4.7249,19.000
451,0.1750,2.0170,1.9197,0.1485,4.7022,4.7156,4.7144,18.137
474,0.1912,2.0234,1.9613,0.1441,4.3818,4.3891,4.3900,18.000
498,0.3993,2.0923,2.0417,0.7672,4.3441,4.3502,4.3502,17.004
523,0.4190,2.0958,2.0222,0.7684,4.2020,4.2019,4.2033,17.000
550,8.3800,2.0577,1.9869,1.0815,4.1264,4.1231,4.1250,16.000
578,8.7688,2.0251,1.9513,1.0648,4.0739,4.0775,4.0786,16.000
607,9.1453,1.9281,1.8733,1.0573,4.0781,4.0737,4.0756,16.000
638,9.5679,2.0155,1.9635,1.0568,4.0322,4.0230,4.0250,16.000
670,10.0017,2.0397,1.9405,1.0459,4.0217,4.0097,4.0129,16.000
704,10.4381,1.9746,1.9274,1.0337,3.8963,3.8886,3.8918,16.000
740,10.9254,2.0096,1.9549,1.0801,3.7433,3.7297,3.7322,16.000
777,11.4206,2.0502,1.9914,1.0828,3.7663,3.7458,3.7456,16.000
816,11.9825,2.0279,1.9531,1.0816,3.5593,3.5352,3.5358,16.000
857,12.5055,2.0541,1.9497,1.0464,3.5825,3.5626,3.5638,16.000
900,13.0573,2.0447,1.9061,1.0644,3.5392,3.5189,3.5224,16.000
945,13.6033,2.1149,1.9401,1.0411,3.4637,3.4331,3.4334,16.000
993,14.1636,2.0835,1.8907,1.0196,3.3627,3.3401,3.3402,16.000
1043,14.7666,2.1819,1.9838,1.0395,3.2955,3.2745,3.2739,16.000
1096,15.4087,2.2158,1.9555,1.0886,3.3190,3.3179,3.3166,16.000
1151,16.0730,2.3093,2.0438,1.1105,3.1324,3.1416,3.1391,16.000
1209,16.7513,2.2840,2.0332,1.0778,3.0787,3.0898,3.0840,16.000
1270,17.4800,2.2443,1.9621,1.0663,3.0185,3.0122,3.0094,16.000
1334,18.2223,2.1815,1.9850,1.0498,3.1611,3.1595,3.1547,16.000
1401,18.9965,2.2603,2.0010,1.0546,3.2113,3.2034,3.2021,16.000
1472,19.7589,2.2849,2.0233,1.0652,3.2082,3.1877,3.1852,16.000
1546,20.5720,2.3194,2.0546,1.0497,3.1207,3.1111,3.1063,16.000
1624,21.3778,2.3898,2.0711,1.0611,2.9350,2.9268,2.9214,16.000
1706,22.2333,2.4213,2.0710,1.0479,2.8731,2.8729,2.8669,16.000
1792,23.1321,2.4664,2.0943,1.0703,2.7623,2.7898,2.7787,16.000
1882,24.0263,2.5293,2.1017,1.0884,2.6701,2.7016,2.6874,16.000
1977,24.9299,2.5721,1.9978,1.0619,2.6707,2.7061,2.6886,16.000
2076,25.9117,2.5192,1.9448,1.0537,2.6957,2.7124,2.6939,16.000
2180,27.5649,2.7549,2.1492,1.4146,2.8811,2.8632,2.8396,16.000
\end{filecontents*}

\begin{abstract}
HyperLogLog and its successors maintain an array of $m = 2^P$
registers, each storing the maximum observed rank, and estimate
cardinality from their harmonic-mean estimator. HLL achieves relative
error $1/\sqrt{m}$ only at high cardinality. At low cardinality
($n \lesssim 2^{P+1}$), the estimate is dominated by sampling noise
and existing systems rely on linear counting or empirically fitted
correction tables. Sparse representations (storing explicit hash
values instead of a fixed register array) extend the near-exact range.
We build on this direction with a three-stage adaptive counter that
progresses from an exact value list to a hash list of distinct
composites, pushing the sparse representation deep into the collision
regime where the raw distinct count is biased low. We correct this
bias by inverting an analytically derived occupancy function via
safeguarded Newton iteration. At precision $P = 10$ and register
width $B = 6$ bits, the corrected hash-list estimate achieves mean
absolute relative error of $0.86\%$ over its regime (compared to
$2.20\%$ for linear counting, $1.98\%$ for register MLE,
and $3.83\%$ for SetSketch with the same bucket count), with zero
error for the value-list portion of the range. The counter converts
to HyperLogLog registers when the hash list saturates. This
three-stage design (value list, hash list, dense representation)
provides a general template for register-based estimators that lack
low cardinality accuracy. The default counter uses fixed-size stack
arrays and is available as an open source
crate.\footnote{\url{https://github.com/LucaCappelletti94/hyperloglog-rs}}
\end{abstract}

\section{Introduction}

Cardinality estimation under space constraints is a foundational
problem in data streaming, database systems, and distributed computing.
The HyperLogLog (HLL) algorithm of Flajolet et al.~\cite{flajolet2007hyperloglog}
hashes each element of a stream and maintains an array of $m = 2^P$
small integer registers, each storing the maximum observed rank (number
of leading zeros plus one) among elements hashing to that bucket. The
cardinality estimate derives from the harmonic-mean estimator and
achieves relative standard error $\sigma \approx 1.04/\sqrt{m}$
in the asymptotic regime ($n \gg m \log m$).

This asymptotic guarantee masks significant error at low
cardinality. Below $n \approx 2^{P+1}$, many registers remain empty and
the harmonic-mean estimator is dominated by the $m$ zero entries,
yielding large relative error. The standard remedy is linear counting
($\hat{n} = -m \ln(Z/m)$ where $Z$ is the number of zero registers),
which is table-free and accurate only in a narrow band near the
transition.

This low-cardinality regime is the operating point of many
real-world workloads. In graph analytics, many
network classes exhibit heavy-tailed degree distributions (power-law or
scale-free graphs) in which a large fraction of nodes have small
degree (often single digits to low hundreds). The HyperBall algorithm
of Boldi and Vigna~\cite{boldi2013hyperball} approximates geometric
centralities (closeness, harmonic centrality) on very large graphs by
maintaining per-node HLL counters during ball expansion. On
heavy-tailed graphs, the per-node HLL counter operates in the low-load
regime for many nodes, where the asymptotic guarantee does
not apply and empirical bias dominates. Our adaptive three-stage
representation addresses this limitation: many nodes remain in
the exact value list or the near-exact hash list, and only
nodes upgrade to the probabilistic register array. This avoids paying
the cost of probabilistic estimation for nodes that do not need it.

HLL++ uses a sparse representation (sorted compressed index-rank
pairs) for very low cardinality, then linear counting in the next
band, then empirically fitted bias-correction tables for the
intermediate regime. Our counter operates in the linear-counting band:
it maintains explicit representations of the inserted data until space
pressure requires compression, progressing through three stages:

\begin{enumerate}
  \item \textbf{Value list.} For the smallest sets, the exact sorted
    list of inserted values is stored with gap compression: the first
    value is encoded absolutely and each subsequent value as the gap
    from its predecessor minus one, using Elias-gamma coding.
    Subtracting one from each gap is always valid (distinct sorted
    values have gaps of at least one), and it makes the common case
    of contiguous values cost a single bit each (gap one encodes as
    zero). The sorted order serves a dual purpose: it enables fast
    set operations (intersection, union) via merge-style traversal,
    and it minimizes the gap codes between consecutive values.
    The effective capacity of the value list depends on the gap
    distribution of the input data: in the best case (contiguous
    values), each element after the first requires a single bit,
    yielding an upper bound equal to the buffer size in bits.
    When the input values follow a global ordering (for example, graph node
    IDs produced by a Layered Label Propagation (LLP)
    reordering~\cite{boldi2010llp} that clusters nodes at multiple
    resolutions and produces an ordering where topologically close
    nodes receive nearby IDs), the gaps are small and capacity
    approaches this bound. Cardinality is exact.
  \item \textbf{Hash list.} When the value list fills, each value is
    replaced by a composite hash: a $w$-bit fingerprint, where $w$ is
    the current composite width, concatenating the $P$-bit bucket
    index, the register rank, and residual hash bits. The sorted list
    of distinct composites is stored with Rice-coded gaps between
    consecutive entries. Cardinality is estimated by inverting an
    analytically derived occupancy function: the expected number of
    distinct composites as a function of cardinality (a birthday-paradox
    problem with non-uniform cell probabilities), inverted via
    safeguarded Newton iteration to recover the cardinality estimate.
  \item \textbf{Dense representation.} When the hash list saturates,
    the counter converts to a representation optimized for large
    cardinality. In the current implementation, this is the standard
    HLL register array.
\end{enumerate}

The novelty of this work lies in the Rice-coded gap packing of the
hash list and the occupancy-inverse bias correction. The packing
pushes cardinality deep enough into the collision regime to make
the correction necessary and impactful. A detailed comparison with
HLL++'s sparse representation follows in \autoref{sec:hllpp}.

The remainder of this paper formalizes the method (\autoref{sec:method}),
derives the occupancy model and its inversion (\autoref{sec:occupancy}),
presents empirical results (\autoref{sec:results}), relates the
approach to HLL++ (\autoref{sec:hllpp}), and discusses
applicability to SetSketch and HyperLogLogLog
(\autoref{sec:discussion}).

\section{Method}
\label{sec:method}

\subsection{Composite Hash Encoding}

Each element $x$ is hashed by a uniform 64-bit hash function $h(x)$.
The top $P$ bits select the register bucket index
$I(x) = \lfloor h(x) / 2^{64-P} \rfloor$. The register rank
$R(x) = \rho(h(x) \bmod 2^{64-P}) + 1$ is the number of leading zeros
of the tail plus one, capped at $2^B - 1$. The composite hash at width
$t = w - P$ tail bits takes three forms depending on $t$ and the
register value:

The composite hash packs three pieces of information into $w$ bits:
the bucket index, the register rank, and extra hash bits for
differentiation. Because the register rank follows a geometric
distribution (small values are common), we store it explicitly only
when it is large, and use the extra hash bits as a proxy when it is
small. This saves space for the common case. Formally:

\begin{definition}[Composite hash forms]
Let $w$ be the current composite width and $t = w - P$ the tail width.
The composite $C(x)$ is a $w$-bit string:
\begin{itemize}
  \item If $t = B$ (smallest width), $C(x)$ concatenates the $P$-bit
    index and the $B$-bit register value. No residual bits remain.
  \item If $t > B$ and the register value $r \leq B + 1$, a flag bit 0
    is followed by $t - 1$ leading hash bits (the most significant
    bits of the original hash tail). The register information is
    implicit: $r$ is small enough that the leading bits carry
    sufficient resolution.
  \item If $t > B$ and $r > B + 1$, a flag bit 1 is followed by the
    $B$-bit register value and $t - 1 - B$ residual hash bits.
\end{itemize}
\end{definition}

The width $w$ starts at a large value (e.g.\ 24 bits for $P = 10$) and
decreases as the list fills. At the maximum width, distinct elements
almost never collide (the composite space is $2^w \gg n$). As $w$
shrinks, the collision rate rises and the number of distinct composites
$D$ begins to undercount $n$. The correction for this undercount is the
occupancy inversion described in \autoref{sec:occupancy}.

\subsection{Rice-Coded Gap Compression}

The composite hashes are maintained in sorted order. Instead of storing
each $w$-bit composite explicitly, we store the gap from each composite
to its predecessor, encoded with a Rice code~\cite{rice1959class}. A Rice code with
parameter $k$ represents an integer $x$ as the quotient
$\lfloor x / 2^k \rfloor$ in unary followed by the $k$-bit remainder.
The expected code length is minimized when $2^k \approx \mu \ln 2$
where $\mu$ is the mean gap.

\begin{lemma}[Optimal Rice parameter]
For $D$ distinct composites uniformly distributed over a $2^w$ space,
the mean gap is $\mu = 2^w / D$. The optimal Rice parameter for a
geometric gap distribution is
\[
k \approx \log_2(\mu \ln 2) = w - \log_2 D - 0.53.
\]
\end{lemma}

\begin{proof}
The gaps between sorted uniform order statistics are geometrically
distributed with mean $\mu = 2^w / D$. Golomb coding is optimal for a
geometric source~\cite{gallager1975golomb}, and Rice coding is Golomb
restricted to a power-of-two parameter $M = 2^k$. The optimal Golomb
parameter for geometric mean $\mu$ is $M \approx \mu \ln 2$, hence
$k = \lfloor \log_2(\mu \ln 2) + 0.5 \rfloor$. Substituting $\mu = 2^w / D$
gives $k \approx w - \log_2 D - 0.53$.
\end{proof}

To obtain a closed form, we substitute the steady-state occupancy
observed empirically: when the width downgrade triggers, the distinct
count clusters near $D \approx 2^{P - 1.53}$ under uniform hashing.
This yields the practical approximation $k \approx w - P + 1$, which
depends only on the current width and precision. The implementation
uses $k = w - P + 1$ for $w > P + B$ and $k = 0$ at the narrowest
width, matching the empirically optimal parameter within one step for
all tested configurations.

The Rice parameter is derived analytically from the encoding structure,
no per-configuration calibration or lookup table is required. The
compression ratio is substantial: at $P = 10$ and $w = 24$, each
composite requires approximately 14 bytes in the value list (Elias-gamma
coded gaps between full 64-bit values spread across the $2^{64}$ domain)
but only 3 to 4 bytes in the sparse representation (Rice-coded gaps
between composites in a space of $2^{24}$).

\subsection{Representation Transitions}

This compression enables the counter to defer register conversion,
but representation changes must be managed. The counter switches
between representations at fixed memory thresholds. The value list holds at most $N_{\text{vl}} \approx 200$
elements for $P = 10$ with uniformly distributed input hashes, but
actual capacity varies with the gap distribution of the encoded data
(higher for clustered or ordered inputs). When a new element
would overflow the value list, all stored values are hashed and the
counter switches to the hash list at the maximum width
$w_{\max}$. The hash list holds elements until the composite
width reaches $w_{\min} = P + B$. At that point, the distinct composites are
materialized into the dense representation. In the current
implementation, the composites are folded into the $m = 2^P$
HLL registers and the counter enters the standard HLL
regime.

The default counter uses fixed-size stack arrays sized at construction
from $(P, B)$ and requires no heap allocation.

The counter supports allocation-free set operations. Merging two value
lists uses a linear two-pointer traversal of the sorted exact values,
producing an exact union. Merging two hash lists merges distinct
composites at the common (coarser) width, re-encodes the result with
Rice-coded gaps, and re-applies the occupancy inversion for the union
estimate. At the register level, the merge reduces to the standard
element-wise maximum of the two HLL arrays.


\section{Occupancy Model and Bias Correction}
\label{sec:occupancy}

\subsection{Expected Distinct Composite Count}

The core of the bias correction is the expected number of distinct
composite values $D$ as a function of cardinality $n$ and width $w$.
This is a birthday-paradox occupancy problem: each of $n$ elements
lands in one of $2^w$ composite ``boxes,'' and $D$ is the number of
non-empty boxes. Because the composite encoding is not uniform
(index bits are uniform, but register bits follow a geometric
distribution and residual bits are uniform), the cell probabilities
are not equal. We group cells by probability and sum over groups.

\begin{theorem}[Occupancy function]
\label{thm:occupancy}
Let $m = 2^P$ and $R = 2^B - 1$ be the maximum register value.
We assume that $h$ is a uniform hash function: each of the $2^{64}$
possible hash values is equally likely for every input element.
The expected distinct composite
count at cardinality $n$ and width $w$ is
\[
\mathbb{E}[D \mid n, w] = \sum_{j} g_j \bigl(1 - (1 - p_j)^n\bigr),
\]
where the sum ranges over cell groups $(g_j, p_j)$ defined by the
composite encoding:
\begin{itemize}
  \item At the smallest width ($t = B$), each non-saturating register
    value $r \in \{1, \dots, R - 1\}$ (where $R - 1 = 2^B - 2$)
    corresponds to $g_r = m$ cells of probability
    $p_r = 2^{-(P+r)}$. The saturating register $R = 2^B - 1$
    corresponds to $g_R = m$ cells of probability
    $p_R = 2^{-(P+2^B-2)}$. Hence
    \begin{equation}
      \mathbb{E}[D \mid n] = m \left(
        \sum_{r=1}^{R-1} \bigl(1 - (1 - 2^{-(P+r)})^n\bigr)
        + \bigl(1 - (1 - 2^{-(P+2^B-2)})^n\bigr)
      \right).
      \label{eq:small}
    \end{equation}
  \item At wider widths ($t > B$), the flag bit partitions the tail
    space. Flag 0 (register $r \leq B + 1$, leading bits stored) has
    $2^{t-2-\ell}$ cells of probability $2^{-(w-1)}$ for each
    $\ell \in \{0, \dots, \min(B, t-2)\}$. Flag 1 (register $r > B +
    1$, register plus residual stored) has $2^{t-1-B}$ cells of
    probability $2^{-(P+r+t-1-B)}$ for each $r \in \{B+2, \dots,
    R-1\}$, with the saturating register at probability
    $2^{-(P+R-1+t-1-B)}$. Hence
    \begin{multline}
      \mathbb{E}[D \mid n, w] = m \biggl(
        \sum_{\ell=0}^{\min(B, t-2)} 2^{t-2-\ell}
        \bigl(1 - (1 - 2^{-(w-1)})^n\bigr) \\
        + \sum_{r=B+2}^{R-1} 2^{t-1-B}
        \bigl(1 - (1 - 2^{-(P+r+t-1-B)})^n\bigr) \\
        + 2^{t-1-B}
        \bigl(1 - (1 - 2^{-(P+R-1+t-1-B)})^n\bigr)
      \biggr).
      \label{eq:wide}
    \end{multline}
\end{itemize}
\end{theorem}

\begin{proof}
Each element independently selects a composite value. The probability
that a given cell with probability $p_j$ is occupied by at least one
of $n$ elements is $1 - (1 - p_j)^n$. Summing over all cells, grouped
by equal probability, gives the result. The cell groups follow from
the composite encoding structure. The index bits are uniform,
contributing the factor $m = 2^P$ to every group. For the tail bits,
the register distribution is geometric: the probability of rank $r$
is $2^{-r}$ for $r < 2^B - 1$ and $2^{-(2^B-2)}$ for the saturating
register. At the smallest width ($t = B$), the tail consists of the
$B$-bit register value, so each register value $r$ corresponds to
exactly $m$ cells (one per index), each with probability
$2^{-(P+r)}$ for $r < 2^B - 1$ and $2^{-(P+2^B-2)}$ for the saturating
register. At wider widths ($t > B$), the flag bit partitions the
tail. Flag 0 stores $t - 1$ leading hash bits for registers
$r \leq B + 1$. For a given number $\ell$ of leading zeros in the
stored hash bits ($\ell \in \{0, \dots, \min(B, t - 2)\}$), there are
$2^{t - 2 - \ell}$ distinct bit patterns, each defining a cell with
probability $2^{-(w - 1)}$. Flag 1 stores the $B$-bit register value
and $t - 1 - B$ residual hash bits for registers $r > B + 1$. Each
such register defines a group of $2^{t - 1 - B}$ cells with
probability $2^{-(P + r + t - 1 - B)}$. The non-uniformity in cell
probabilities arises from the geometric register distribution and the
flag-bit partition at wide widths.
\end{proof}

\subsection{Occupancy Inversion}

The estimator $\hat{n}$ is the inverse of the occupancy function at the
observed distinct count $D$:
\begin{equation}
  \hat{n} = g^{-1}(D) \quad \text{where} \quad g(n) = \mathbb{E}[D \mid n, w].
  \label{eq:inverse}
\end{equation}

We solve $g(n) = D$ by safeguarded Newton iteration. The derivative
required for the Newton step is
\begin{equation}
  g'(n) = \sum_{j} g_j \bigl(-\ln(1 - p_j)\bigr) (1 - p_j)^n,
  \label{eq:deriv}
\end{equation}
which is always positive (the occupancy function is strictly
increasing), ensuring convergence. The safeguard maintains a bracket
$[\text{low}, \text{high}]$ initialized at $[D, D]$ and expanded until
$g(\text{high}) \geq D$. Each Newton step $n_{\text{next}} = n -
(g(n) - D) / g'(n)$ is accepted only if it falls within the bracket,
otherwise, a bisection step is taken. Convergence is declared when the
relative step falls below $10^{-12}$.

Because $D$ is a random variable and the inversion $g^{-1}$ is
nonlinear, Jensen's inequality implies that the estimator carries a
small bias: $\mathbb{E}[g^{-1}(D)] \neq g^{-1}(\mathbb{E}[D]) = n$.
The empirical MRE values in \autoref{tab:error} (all near zero)
confirm that this bias is negligible in practice for the
configurations tested.

The predicted variance of the estimator uses a delta-method
approximation that treats the multinomial cell indicators as if they
were independent. For the non-uniform cell probabilities of the
composite encoding, this is an approximation: the exact variance
requires cross terms between cells of different probability. The
approximation is accurate to within a few percent for the
configurations tested, as confirmed by comparing predicted and
empirical standard errors.

\begin{algorithm}
\caption{Occupancy-inversion estimator}
\label{alg:occupancy}
\begin{algorithmic}[1]
\Require Observed distinct count $D$, current width $w$
\Ensure Cardinality estimate $\hat{n}$
\State Compute cell groups $(g_j, p_j)$ from encoding at width $w$
\Function{Occupancy}{$n$}
  \State $E \gets \sum_j g_j \bigl(1 - (1 - p_j)^n\bigr)$
  \State \Return $E$
\EndFunction
\Function{OccupancyDerivative}{$n$}
  \State $E' \gets \sum_j g_j \bigl(-\ln(1 - p_j)\bigr) (1 - p_j)^n$
  \State \Return $E'$
\EndFunction
\State $\text{low} \gets D$, $\text{high} \gets D$
\While{$\text{Occupancy}(\text{high}) < D$}
  \State $\text{high} \gets 2 \cdot \text{high}$
\EndWhile
\State $n \gets (\text{low} + \text{high}) / 2$
\Repeat
  \If{$\text{Occupancy}(n) > D$}
    \State $\text{high} \gets n$
  \Else
    \State $\text{low} \gets n$
  \EndIf
  \State $c \gets n - (\text{Occupancy}(n) - D) / \text{OccupancyDerivative}(n)$
  \If{$\text{low} < c < \text{high}$}
    \State $n_{\text{next}} \gets c$
  \Else
    \State $n_{\text{next}} \gets (\text{low} + \text{high}) / 2$
  \EndIf
\Until{$|n_{\text{next}} - n| / n \leq 10^{-12}$}
\State \Return $n_{\text{next}}$
\end{algorithmic}
\end{algorithm}

The computational cost is $O(2^B)$ per estimate: the sums in
\eqref{eq:small} and \eqref{eq:wide} have $O(2^B)$ terms, and the
Newton loop converges in a bounded number of iterations (typically 5 to 10).
The cost depends on neither $n$ nor $m = 2^P$, and the stored entries
are never scanned.

\section{Empirical Results}
\label{sec:results}

We evaluate the method at precision $P = 10$ ($m = 1024$ registers) and
register width $B = 6$ ($R = 63$). The sparse regime spans
cardinalities from approximately $n = 52$ (end of the value list) to
$n \approx 2153$ (conversion to registers). This range falls within the
linear-counting regime of the register estimator (threshold $n = 2320$
for this configuration), and the HLL++ empirical bias-correction tables
do not engage until $n > 2320$. For each cardinality, we report the
mean absolute relative error over 256 random seeds.

\begin{table}[t]
\centering
\caption{Aggregate error over the sparse regime ($P = 10$, $B = 6$,
256 seeds). MAE is mean absolute error in elements, MRE is mean
signed relative error, MSE is mean squared error, and MARE is mean
absolute relative error (cardinality-weighted over the sparse regime).
The baseline register estimate uses linear counting in this range.}
\label{tab:error}
\begin{tabular}{lrrrr}
\toprule
Estimator & MAE & MRE & MSE & MARE \\
\midrule
Uncorrected ($D$) & 208.1 & $-13.98\%$ & 75551 & 13.98\% \\
Registers (linear counting) & 25.8 & $+0.17\%$ & 1432 & 2.20\% \\
Registers (MLE) & 22.4 & $+0.12\%$ & 1052 & 1.98\% \\
Corrected (occupancy inverse) & 11.3 & $-0.02\%$ & 292 & 0.85\% \\
SetSketch (b=2) & 35.4 & $-0.01\%$ & 2332 & 3.83\% \\
SetSketch (b=1.2) & 35.4 & $-0.02\%$ & 2324 & 3.83\% \\
\bottomrule
\end{tabular}
\end{table}

The corrected estimate is $2.3\times$ better than linear counting on
MAE and $4.9\times$ on MSE. The register MLE estimator of
Ertl~\cite{ertl2017probabilistic} improves over linear counting
(MAE 22.4 vs 25.8, MARE 1.98\% vs 2.20\%) but remains $2.3\times$
worse than the corrected hash list on MARE (1.98\% vs 0.86\%).
SetSketch with the same bucket count ($m = 1024$) achieves MAE of
35.4 across all tested $b$ values, roughly $1.4\times$ worse than
linear counting and $4.5\times$ worse than the corrected estimate.
All three correction methods (MLE, occupancy inverse, and SetSketch)
are effectively unbiased (MRE near zero), while the uncorrected
distinct count is biased low by 14\% due to collisions.

The critical distinction between the hash list and
SetSketch appears at low cardinality. The three-stage design
provides exact cardinality through the value-list stage ($n \leq 52$)
and near-exact estimates in the early hash-list regime (MARE below
0.01\%). SetSketch has no exact-counting mode: it is purely
stochastic from $n = 1$, producing MARE of 69\% at $n = 1$,
and only reaching the 3\% range beyond $n = 500$. This structural
gap means SetSketch cannot replace a multi-stage adaptive
representation when low-cardinality accuracy matters.

\begin{figure}[t]
\centering
\resizebox{\linewidth}{!}{\begin{tikzpicture}
\begin{axis}[
  width=0.90\linewidth, height=6cm, scale only axis,
  axis y line*=left,
  xlabel={true cardinality $n$},
  ylabel={mean absolute relative error (\%)},
  xmin=0, xmax=2200, ymode=log, ymin=0.001, ymax=100,
  legend cell align=left, legend columns=3,
  legend style={font=\scriptsize, at={(0.5,-0.20)}, anchor=north,
    /tikz/every even column/.append style={column sep=0.5cm}},
  grid=both, grid style={gray!20},
  tick label style={font=\footnotesize},
  label style={font=\footnotesize},
]
\draw[gray, densely dashed] (axis cs:52,0.001) -- (axis cs:52,100);
\draw[gray, densely dashed] (axis cs:2153,0.001) -- (axis cs:2153,100);
\node[gray, font=\scriptsize, rotate=90, anchor=south]
  at (axis cs:52,0.09) {value list ends};
\node[gray, font=\scriptsize, rotate=90, anchor=south]
  at (axis cs:2153,0.09) {registers begin};
\addplot[red, thick] table[col sep=comma, x=n, y=uncorrected]
  {hll_plot_data.csv};
\addlegendentry{uncorrected ($D$)}
\addplot[orange, thick] table[col sep=comma, x=n, y=hllpp]
  {hll_plot_data.csv};
\addlegendentry{registers (linear counting)}
\addplot[green!60!black, thick, dashed] table[col sep=comma, x=n, y=hll_mle]
  {hll_plot_data.csv};
\addlegendentry{registers (MLE)}
\addplot[blue, thick] table[col sep=comma, x=n, y=corrected]
  {hll_plot_data.csv};
\addlegendentry{corrected (occupancy inverse)}
\addplot[purple, thick] table[col sep=comma, x=n, y=setsketch_b2]
  {hll_plot_data.csv};
\addlegendentry{SetSketch (b=2)}
\addplot[cyan, thick, dashed] table[col sep=comma, x=n, y=setsketch_b1p2]
  {hll_plot_data.csv};
\addlegendentry{SetSketch (b=1.2)}
\addlegendimage{black, densely dotted, thick}
\addlegendentry{avg composite width (right axis)}
\end{axis}
\begin{axis}[
  width=0.90\linewidth, height=6cm, scale only axis,
  axis y line*=right, axis x line=none,
  xmin=0, xmax=2200, ymin=14, ymax=26,
  ylabel={avg width (bits)},
  tick label style={font=\footnotesize},
  label style={font=\footnotesize},
]
\addplot[black, densely dotted, thick, restrict y to domain=1:100,
  unbounded coords=discard]
  table[col sep=comma, x=n, y=bits] {hll_plot_data.csv};
\end{axis}
\end{tikzpicture}}
\caption{Mean absolute relative error over 256 seeds ($P = 10$, $B =
6$). The uncorrected count $D$ is biased low. The register baseline
uses linear counting. The occupancy inverse stays below both.
SetSketch (same $m = 1024$) has higher error and lacks an
exact-counting mode. The dotted curve (right axis) is the average
composite width, falling from 24 to 16 bits. Dashed verticals mark
representation transitions.}
\label{fig:error}
\end{figure}

\section{Relation to HyperLogLog++}
\label{sec:hllpp}

HLL++~\cite{heule2013hyperloglog} (deployed in Google's infrastructure)
introduces three improvements over the original HLL: a 64-bit hash
function, a sparse representation for low cardinality, and an
empirically fitted bias-correction table. The HLL++ sparse
representation and our compressed sparse representation share the
same fundamental idea: defer the register array by storing explicit
hash data in a compressed sorted list. HLL++ stores index-rank
pairs (one per bucket, residuals discarded); the hash list stores
composites (index plus rank plus residuals, multiple per bucket).

\paragraph{What HLL++ stores.}
The HLL++ sparse representation keeps one entry per occupied bucket:
a pair of bucket index and rank $(\text{idx}, \rho(w))$ at elevated
precision ($p' = 25$). If two elements hash to the same bucket, only
the one with the higher rank survives. The residual hash bits beyond
the index and rank are discarded. The sorted list is compressed with
variable-length integer encoding and difference coding, using roughly
8 bits per entry at $P = 14$. When the compressed list exceeds $6m$
bits in size, HLL++ converts to the dense register array. At query time, if
the representation is still sparse, linear counting on the number of
distinct indices gives the estimate.

\paragraph{What the hash list stores.}
The hash list keeps distinct \emph{composites}: each
entry concatenates the $P$-bit bucket index, the register rank,
\emph{and residual hash bits} from the original hash. Two elements
with different ranks both survive as separate composites. The sorted
distinct composites are packed with Rice-coded gaps between
consecutive entries. The composite width $w$ starts large (e.g.\ 24
bits at $P = 10$) and shrinks as the list fills, making the gaps
smaller and the Rice codes shorter.

\paragraph{Why retaining residual bits matters.}
By keeping residual hash bits, the hash list carries more
cardinality per byte than HLL++'s sparse representation. The Rice-coded gap
compression makes this affordable: at $P = 10$ and $w = 24$, each
composite costs 3-4 bytes. The wider composite space ($2^{24}$ versus
$2^{14}$ buckets) is offset by Rice coding over small gaps. The extra
information is what pushes the representation much further: the hash
list carries cardinality deep into the collision regime
(at $P = 10$), where HLL++ would have already converted to registers
($n \approx 534$ at $P = 10$, Heule et al.~\cite{heule2013hyperloglog}).

\paragraph{Why pushing further requires bias correction.}
Because the hash list reaches cardinalities where $n$
approaches the composite space $2^w$, collisions between distinct elements
producing the same composite become frequent. The raw distinct count
$D$ underestimates $n$ by over $14\%$ near conversion. The
occupancy-inverse correction is what makes this regime usable: it
models the non-uniform collision structure of the composite encoding
and inverts the expected distinct count to recover $n$. HLL++ does not
face this problem because its sparse representation converts to
registers before collisions become significant, relying instead on
linear counting (which is accurate in the narrow band where HLL++
operates) and its empirical bias tables for the intermediate regime.

\paragraph{Complementarity.}
The hash list and HLL++ are complementary. The hash
list covers the linear-counting regime of the register
estimator ($n \lesssim 2153$ at $P = 10$), where HLL++ relies on
linear counting and empirical bias tables. The hash list
provides a table-free alternative with lower error across the entire regime.

\section{Discussion: SetSketch and HyperLogLogLog}
\label{sec:discussion}

SetSketch~\cite{ertl2021setsketch} generalizes both HLL and MinHash
by storing per-bucket maximums of exponentially distributed hash
values instead of uniform ranks. A continuous parameter $b$
interpolates between the HLL limit ($b = 2$, HLL-equivalent register
statistics) and the MinHash limit ($b \to 1$, equivalent to MinHash
via a monotone transformation). Its cardinality estimator uses a
closed-form harmonic-mean variant with analytically derived
corrections for small and large cardinalities, requiring no empirical
calibration.
We evaluate SetSketch across three $b$ values: $b = 2$ (HLL limit),
$b = 1.2$, and $b = 1.001$ (near the MinHash limit).

Our empirical comparison (\autoref{fig:error}, \autoref{tab:error})
evaluates SetSketch with the same bucket count ($m = 1024$)
against the hash list. In the sparse regime SetSketch achieves MARE
of 3.83\% across all tested $b$ values, higher than both linear
counting (2.20\%), the register MLE (1.98\%), and the occupancy
inverse (0.86\%). The different $b$ values produce nearly identical
results, confirming the estimator is insensitive to this parameter
in the sparse regime. The dominant limitation is low-cardinality
accuracy: SetSketch lacks an exact-counting mode and is purely
stochastic from $n = 1$, yielding MARE of 69\% at $n = 1$ and
12\% at $n = 50$. The three-stage design avoids this entirely through
the value-list stage (exact through $n = 52$) and the near-exact
early hash-list regime.

The occupancy model of the hash list is not directly applicable
to SetSketch because its register values follow an exponential
distribution rather than the non-uniform composite encoding
structure. An adapted model would need a separate derivation for the
order-statistic distribution of per-bucket maxima under exponential
hashing.

\subsection{HyperLogLogLog}
HyperLogLogLog~\cite{karppa2022hyperlogloglog} compresses the
HyperLogLog register array by storing a base value $B$ and per-register
offsets $M[j] - B$ in a dense array, with overflowing registers kept
sparsely as $(j, M[j])$ pairs. This reduces the sketch from
$O(m \log \log n)$ bits to $m \log_2 \log_2 \log_2 n + O(m + \log \log n)$
bits while preserving all estimation properties and mergeability in
the compressed domain.

The hash list and HyperLogLogLog address different
regimes: HyperLogLogLog compresses the register representation at
high cardinality, while the hash list replaces registers
entirely at low cardinality. The two techniques are complementary:
the hash list defers the register abstraction, and
HyperLogLogLog compresses it when it is eventually needed.

\section{Conclusion}

We presented a three-stage adaptive representation for cardinality
estimation that achieves sub-register error in the low-cardinality
regime. The hash list with Rice-coded gaps and
bias correction reduces mean absolute relative error by $2.6\times$
over linear counting, $2.3\times$ over register MLE, and
$4.5\times$ over SetSketch. The MLE approach of
Ertl~\cite{ertl2017probabilistic} narrows the gap between linear
counting and the corrected estimate (MARE 1.98\% vs 2.20\% vs
0.86\%) but cannot overcome the fundamental limitation: at low
cardinality the register statistics are dominated by zeros, and no
post-hoc correction on the register array alone can recover the lost
information. The three-stage design (value list, hash list, then dense
representation) avoids this structural gap entirely by keeping data
explicit until compression is unavoidable. The counter converts to HLL
registers when the hash list saturates, so it drops directly into
existing HLL++ pipelines. The $O(d)$ insert cost in
the hash list is the primary limitation, where $d$ is the number of
stored distinct composites: each insertion requires a sorted splice
into the gap-encoded bitstream, which involves decompressing,
inserting, and recompressing a portion of the list. Because the
Rice-coded packing efficiency depends on the gap distribution of the
input hashes, the estimator variance is data-dependent rather than a
function of cardinality alone. Future work includes deriving an
occupancy model for SetSketch's exponential register distribution and
formalizing the occupancy theorem in a proof assistant (e.g.\ Lean)
to guarantee correctness of the cell probability assignments. The
implementation is available as an open source
crate.\footnote{\url{https://github.com/LucaCappelletti94/hyperloglog-rs}}

\begin{thebibliography}{10}

\bibitem{flajolet2007hyperloglog}
Philippe Flajolet, Eric Fusy, Olivier Gandouet, and Frederic Meunier.
\newblock HyperLogLog: the analysis of a near-optimal cardinality
  estimation algorithm.
\newblock In \emph{Proceedings of the 2007 conference on analysis of
  algorithms}, pages 127--146, 2007.

\bibitem{heule2013hyperloglog}
Stefan Heule, Marc Nunkesser, and Alexander Hall.
\newblock HyperLogLog in practice: algorithmic engineering of a state of the
  art cardinality estimation algorithm.
\newblock In \emph{EDBT/ICDT}, pages 3--14, 2013.


\bibitem{ertl2021setsketch}
Otmar Ertl.
\newblock SetSketch: filling the gap between MinHash and HyperLogLog.
\newblock \emph{Proceedings of the VLDB Endowment}, 14(11):2244--2257,
  2021.

\bibitem{ertl2017probabilistic}
Otmar Ertl.
\newblock Probabilistic analysis of the HyperLogLog algorithm.
\newblock \emph{Algorithmica}, 79(2):519--535, 2017.

\bibitem{karppa2022hyperlogloglog}
Matti Karppa and Rasmus Pagh.
\newblock HyperLogLogLog: cardinality estimation with one log more.
\newblock \emph{arXiv preprint arXiv:2205.11327}, 2022.

\bibitem{gallager1975golomb}
Robert~G. Gallager and Philip~C. van Voorhis.
\newblock Optimal noninteger expansions of digital data.
\newblock \emph{Bell System Technical Journal}, 54(8):1415--1423, 1975.

\bibitem{rice1959class}
Solomon~W. Rice.
\newblock Class of two-dimensional transform coding techniques.
\newblock \emph{Bell System Technical Journal}, 38(3):411--432, 1959.
\bibitem{boldi2010llp}
Paolo Boldi, Marco Rosa, Matteo Santini, and Sebastiano Vigna.
\newblock Layered label propagation: a unified approach to hierarchy
  extraction and ordering in graphs.
\newblock \emph{arXiv preprint arXiv:1011.5425}, 2010.
\bibitem{boldi2013hyperball}
Paolo Boldi and Sebastiano Vigna.
\newblock In-core computation of geometric centralities with HyperBall:
  a hundred billion nodes and beyond.
\newblock \emph{arXiv preprint arXiv:1308.2144}, 2013.


\end{thebibliography}

\end{document}
